% GENERATED FILE. Edit canonical lessons, then run npm run generate:books.
% canonical-source-hash: fnv1a64:1c7473a70353ceed
% canonical-lessons: IT-C476-deludere, IT-C476-delusione, IT-C476-consolare, IT-C476-pianto, IT-C476-singhiozzare

\chapter{Una grande delusione}
\label{ch:it-una-grande-delusione}
\hlchaptermodality{476}

\begin{chapteropening}
\textbf{By the end of this chapter:} \emph{I can talk about a disappointment and how friends comfort me.}

Complete the last lesson of chapter 476: I can talk about a disappointment and how friends comfort me.
\end{chapteropening}

\section[deludere]{deludere — to disappoint, to let down}

\label{lesson:IT-C476-deludere}



\begin{quote}
Before the new one: say the Italian for relief, then the Italian for to calm down.
\end{quote}

\subsection*{You'll want to know: deludere}
\textbf{deludere} — \textquotedblleft{}to disappoint, to let down\textquotedblright{}.

\begin{quote}
Il film mi ha deluso: mi aspettavo di più.
\end{quote}

\begin{quote}
\emph{The film disappointed me: I expected more.}
\end{quote}

A false friend: it means to disappoint, not to delude. The past participle is \emph{deluso}.

\begin{cousinweb}
\emph{deludere} is a learned borrowing from Latin \emph{dēlūdere} \textquotedblleft{}to make fun of, to deceive\textquotedblright{} (Lewis \& Short), built on \emph{lūdere} \textquotedblleft{}to play\textquotedblright{}.
\end{cousinweb}

\begin{grammarlens}[title={What you've built}]
One more word for telling, explaining and describing at B1.
\end{grammarlens}

\subsection*{Your turn}
\begin{itemize}
\raggedright
  \item \emph{Say it:} \emph{deludere}
  \item \emph{Say it:} \emph{deludere}, once more
  \item \emph{Say it:} a sentence of your own with \emph{deludere}
  \item \emph{Recall:} say the Italian for relief, then the Italian for to calm down, then say \emph{deludere} again
\end{itemize}

\begin{tcolorbox}[breakable,colback=teal!4,colframe=teal!35!black,title={Before you move on}]
What does deludere mean? (\textquotedblleft{}To disappoint, to let down\textquotedblright{}.) Say it once more.
\end{tcolorbox}
\section[la delusione]{la delusione — disappointment}

\label{lesson:IT-C476-delusione}



\begin{quote}
Before the new one: say the Italian for to calm down, then the Italian for to disappoint.
\end{quote}

\subsection*{You'll want to know: la delusione}
\textbf{la delusione} — \textquotedblleft{}disappointment\textquotedblright{}.

\begin{quote}
Non passare l'esame è stata una grande delusione.
\end{quote}

\begin{quote}
\emph{Failing the exam was a big disappointment.}
\end{quote}

\emph{Una delusione d'amore} — a heartbreak.

\begin{cousinweb}
\emph{delusione} is a learned borrowing from Late Latin \emph{dēlūsiō} \textquotedblleft{}a deceiving\textquotedblright{} (Lewis \& Short), from \emph{dēlūdere}; in Italian it came to mean the feeling when hopes are let down.
\end{cousinweb}

\begin{grammarlens}[title={What you've built}]
One more word for telling, explaining and describing at B1.
\end{grammarlens}

\subsection*{Your turn}
\begin{itemize}
\raggedright
  \item \emph{Say it:} \emph{la delusione}
  \item \emph{Say it:} \emph{la delusione}, once more
  \item \emph{Say it:} a sentence of your own with \emph{la delusione}
  \item \emph{Recall:} say the Italian for to calm down, then the Italian for to disappoint, then say \emph{la delusione} again
\end{itemize}

\begin{tcolorbox}[breakable,colback=teal!4,colframe=teal!35!black,title={Before you move on}]
What does la delusione mean? (\textquotedblleft{}Disappointment\textquotedblright{}.) Say it once more.
\end{tcolorbox}
\section[consolare]{consolare — to comfort, to console}

\label{lesson:IT-C476-consolare}



\begin{quote}
Before the new one: say the Italian for to disappoint, then the Italian for disappointment.
\end{quote}

\subsection*{You'll want to know: consolare}
\textbf{consolare} — \textquotedblleft{}to comfort, to console\textquotedblright{}.

\begin{quote}
Ho consolato la mia amica dopo la brutta notizia.
\end{quote}

\begin{quote}
\emph{I comforted my friend after the bad news.}
\end{quote}

Reflexive \emph{consolarsi} — to cheer oneself up.

\begin{cousinweb}
\emph{consolare} is a learned borrowing from Latin \emph{cōnsōlārī} \textquotedblleft{}to comfort, to cheer\textquotedblright{} (EtymDB; Lewis \& Short).
\end{cousinweb}

\begin{grammarlens}[title={What you've built}]
One more word for telling, explaining and describing at B1.
\end{grammarlens}

\subsection*{Your turn}
\begin{itemize}
\raggedright
  \item \emph{Say it:} \emph{consolare}
  \item \emph{Say it:} \emph{consolare}, once more
  \item \emph{Say it:} a sentence of your own with \emph{consolare}
  \item \emph{Recall:} say the Italian for to disappoint, then the Italian for disappointment, then say \emph{consolare} again
\end{itemize}

\begin{tcolorbox}[breakable,colback=teal!4,colframe=teal!35!black,title={Before you move on}]
What does consolare mean? (\textquotedblleft{}To comfort, to console\textquotedblright{}.) Say it once more.
\end{tcolorbox}
\section[il pianto]{il pianto — crying, tears}

\label{lesson:IT-C476-pianto}



\begin{quote}
Before the new one: say the Italian for disappointment, then the Italian for to comfort.
\end{quote}

\subsection*{You'll want to know: il pianto}
\textbf{il pianto} — \textquotedblleft{}crying, tears\textquotedblright{}.

\begin{quote}
Dalla stanza si sentiva il pianto di un bambino.
\end{quote}

\begin{quote}
\emph{From the room you could hear a child crying.}
\end{quote}

The noun for \emph{piangere}; \emph{scoppiare in un pianto} — to burst into tears.

\begin{cousinweb}
\emph{pianto} is inherited from Latin \emph{planctus} \textquotedblleft{}a beating of the breast in grief, a lament\textquotedblright{} (EtymDB; Lewis \& Short), from \emph{plangere}.
\end{cousinweb}

\begin{grammarlens}[title={What you've built}]
One more word for telling, explaining and describing at B1.
\end{grammarlens}

\subsection*{Your turn}
\begin{itemize}
\raggedright
  \item \emph{Say it:} \emph{il pianto}
  \item \emph{Say it:} \emph{il pianto}, once more
  \item \emph{Say it:} a sentence of your own with \emph{il pianto}
  \item \emph{Recall:} say the Italian for disappointment, then the Italian for to comfort, then say \emph{il pianto} again
\end{itemize}

\begin{tcolorbox}[breakable,colback=teal!4,colframe=teal!35!black,title={Before you move on}]
What does il pianto mean? (\textquotedblleft{}Crying, tears\textquotedblright{}.) Say it once more.
\end{tcolorbox}
\section[singhiozzare]{singhiozzare — to sob; to have hiccups}

\label{lesson:IT-C476-singhiozzare}



\begin{quote}
Before the new one: say the Italian for to comfort, then the Italian for crying.
\end{quote}

\subsection*{You'll want to know: singhiozzare}
\textbf{singhiozzare} — \textquotedblleft{}to sob; to have hiccups\textquotedblright{}.

\begin{quote}
Il bambino singhiozzava e non riusciva a parlare.
\end{quote}

\begin{quote}
\emph{The child was sobbing and could not speak.}
\end{quote}

\emph{Il singhiozzo} is a sob, and also the hiccups: \emph{ho il singhiozzo}.

\begin{cousinweb}
\emph{singhiozzare} is inherited from Vulgar Latin \emph{singlutiāre} (EtymDB), from Latin \emph{singultus} \textquotedblleft{}a sob, a hiccup\textquotedblright{} (Lewis \& Short).
\end{cousinweb}

\begin{grammarlens}[title={What you've built}]
One more word for telling, explaining and describing at B1.
\end{grammarlens}

\subsection*{Your turn}
\begin{itemize}
\raggedright
  \item \emph{Say it:} \emph{singhiozzare}
  \item \emph{Say it:} \emph{singhiozzare}, once more
  \item \emph{Say it:} a sentence of your own with \emph{singhiozzare}
  \item \emph{Recall:} say the Italian for to comfort, then the Italian for crying, then say \emph{singhiozzare} again
\end{itemize}

\begin{tcolorbox}[breakable,colback=teal!4,colframe=teal!35!black,title={Before you move on}]
What does singhiozzare mean? (\textquotedblleft{}To sob; to have hiccups\textquotedblright{}.) Say it once more.
\end{tcolorbox}
